% Author: SuFan; MIT License. Illustrative signed values, not measured data.
\documentclass{article}
\usepackage[margin=20mm]{geometry}
\usepackage[libraries=none]{gradbars}
\pagestyle{empty}
\begin{document}
\section*{Positive and negative colors on one track}
\gradbarssetup{min=-50,max=50,error=8,width=72mm,height=6pt,
  precision=1,label width=3em}
\renewcommand{\arraystretch}{1.8}
\newcommand{\themerow}[1]{#1 & $-24.0$ & \gradsignedbar[theme=#1,label=none]{-24}{24} & $24.0$\\}
\begin{tabular}{lrcl}
Theme & Negative & Shared scale: $[-50,50]$ & Positive\\\hline
\themerow{blue}
\themerow{teal}
\themerow{solid}
\themerow{gray}
\themerow{lbyellow}
\themerow{viblue}
\themerow{cyblu}
\themerow{violet}
\themerow{orange}
\themerow{forest}
\themerow{accessible}
\end{tabular}
\par\medskip
Each explicit theme sets the negative fill to the RGB complement of its
positive endpoints. Both values share one track and the central zero line:
the negative value extends left and the positive value extends right.
Whiskers show a declared symmetric error of 8 for each value.
The built-in \texttt{\string\gradsignedbar} command draws both observations;
its two arguments are the nonpositive and nonnegative values, respectively.
\par\medskip
Override after the theme for a custom solid negative fill:
\par\smallskip
\gradsignedbar[theme=blue,negative fill=solid,negative color=lightred]{-24}{24}
\end{document}
